\chapter{Capstone project}
\label{ch:capstone}

\begin{quote}
\emph{``The best way to learn data analysis is to analyze data. Not toy data --- real data, with real messiness, about real health questions.''}
\end{quote}

% ============================================================
\section{Project guidelines}

The capstone project is your opportunity to apply every skill from this course to a real health question. You will work individually or in pairs.

\subsection{Timeline}

\begin{itemize}
  \item \textbf{Week 1.} Choose a project, download the data, perform initial exploration.
  \item \textbf{Week 2.} Complete the analysis, generate figures and tables, draft the report.
  \item \textbf{Week 3.} Finalize the report and prepare a 5-minute oral presentation.
\end{itemize}

\subsection{Expectations}

Your project must include:
\begin{enumerate}
  \item \textbf{A clear health question.} Not ``explore the data'' but ``is diabetes prevalence associated with BMI after adjusting for age and sex?''
  \item \textbf{Real public data.} No simulated datasets. All data sources must be cited.
  \item \textbf{Reproducible code.} A Jupyter notebook that runs from top to bottom without errors.
  \item \textbf{At least 3 visualizations.} Relevant, labeled, and interpretable.
  \item \textbf{At least 1 statistical test or model.} With proper interpretation.
  \item \textbf{A written report.} 5--8 pages following the template below.
  \item \textbf{An oral presentation.} 5 minutes with slides.
\end{enumerate}

\begin{warningbox}
Academic integrity: you may use AI tools (ChatGPT, GitHub Copilot) to help write code, but you must understand every line. During the oral presentation, you will be asked to explain your code and interpret your results. Copy-pasting code you do not understand will be obvious.
\end{warningbox}

% ============================================================
\section{Suggested projects}

Choose one of the five projects below, or propose your own (subject to instructor approval). Each project specifies the objective, dataset(s), required analyses, and expected deliverables.

% ------------------------------------------------------------
\subsection{Project 1: Diabetes risk factor analysis}

\textbf{Objective.} Identify the strongest risk factors for type 2 diabetes and build a predictive model.

\textbf{Datasets.}
\begin{itemize}
  \item \textbf{Pima Indians Diabetes Dataset} (UCI/Kaggle):\\
  \url{https://raw.githubusercontent.com/jbrownlee/Datasets/master/pima-indians-diabetes.data.csv}\\
  768 female patients, 8 features (pregnancies, glucose, blood pressure, skin thickness, insulin, BMI, diabetes pedigree function, age), binary outcome.
  \item \textbf{CDC BRFSS} (Behavioral Risk Factor Surveillance System):\\
  \url{https://www.cdc.gov/brfss/annual_data/annual_data.htm}\\
  Annual survey of 400\,000+ US adults. Download the most recent year's CSV.
\end{itemize}

\textbf{Required analyses.}
\begin{enumerate}
  \item Load and clean the Pima dataset. Handle zeros in glucose, blood pressure, BMI (these are missing values coded as 0).
  \item Exploratory analysis: distributions, correlations, box plots by diabetes status.
  \item Logistic regression: which variables are significant predictors? Report odds ratios with 95\% confidence intervals.
  \item Compare logistic regression with a random forest classifier. Report accuracy, sensitivity, specificity, and AUC-ROC.
  \item (Bonus) Repeat the risk factor analysis on BRFSS data and compare findings.
\end{enumerate}

\textbf{Deliverables.}
\begin{itemize}
  \item Correlation heatmap of all features.
  \item ROC curves comparing logistic regression and random forest.
  \item Table of odds ratios with confidence intervals.
  \item Written interpretation: which risk factors matter most clinically?
\end{itemize}

\begin{clinicalnote}
The Pima Indians dataset is historically important but has ethical considerations. The data was collected from the Akimel O'odham (Pima) community, which has one of the highest rates of type 2 diabetes in the world. When presenting, acknowledge the population and avoid generalizing findings to other groups.
\end{clinicalnote}

% ------------------------------------------------------------
\subsection{Project 2: COVID-19 impact analysis across African countries}

\textbf{Objective.} Analyze how COVID-19 affected African countries differently and explore correlations with health system capacity and socioeconomic factors.

\textbf{Datasets.}
\begin{itemize}
  \item \textbf{JHU CSSE COVID-19 data}:\\
  \url{https://github.com/CSSEGISandData/COVID-19/tree/master/csse_covid_19_data/csse_covid_19_time_series}\\
  Global confirmed cases, deaths, and recovered (daily time series).
  \item \textbf{Our World in Data} (vaccination, testing):\\
  \url{https://raw.githubusercontent.com/owid/covid-19-data/master/public/data/owid-covid-data.csv}\\
  Comprehensive dataset with cases, deaths, vaccinations, testing, hospital data, and demographic indicators.
  \item \textbf{World Bank Open Data} (GDP, health expenditure, hospital beds):\\
  \url{https://data.worldbank.org/}\\
  Use indicator codes: \texttt{SH.XPD.CHEX.PC.CD} (health expenditure per capita), \texttt{SH.MED.BEDS.ZS} (hospital beds per 1\,000), \texttt{NY.GDP.PCAP.CD} (GDP per capita).
  \item \textbf{Africa CDC}:\\
  \url{https://africacdc.org/covid-19/}\\
  Continental-level dashboards and reports.
\end{itemize}

\textbf{Required analyses.}
\begin{enumerate}
  \item Select 10--15 African countries across different subregions (West, East, Southern, North, Central Africa).
  \item Plot epidemic curves (7-day rolling average) for each country. Identify wave patterns.
  \item Compute key metrics: total cases per million, deaths per million, case fatality rate, vaccination rate.
  \item Merge with World Bank indicators. Compute correlations between COVID-19 outcomes and health system capacity (health expenditure, hospital beds, physician density).
  \item Build a choropleth map of case fatality rates across Africa.
  \item (Bonus) Run a multiple regression: does health expenditure per capita predict case fatality rate after controlling for median age and GDP?
\end{enumerate}

\textbf{Deliverables.}
\begin{itemize}
  \item Epidemic curves for all selected countries (small multiples or grouped plot).
  \item Scatter plot of health expenditure vs.\ case fatality rate.
  \item Choropleth map of Africa showing COVID-19 burden.
  \item Summary table with key metrics for all countries.
\end{itemize}

% ------------------------------------------------------------
\subsection{Project 3: Maternal mortality determinants}

\textbf{Objective.} Investigate which health system and socioeconomic factors are associated with maternal mortality across countries.

\textbf{Datasets.}
\begin{itemize}
  \item \textbf{WHO Global Health Observatory} --- Maternal mortality ratio (MMR):\\
  \url{https://apps.who.int/gho/athena/api/GHO/MDG_0000000026?format=csv}
  \item \textbf{WHO GHO} --- Skilled birth attendance:\\
  \url{https://apps.who.int/gho/athena/api/GHO/MDG_0000000025?format=csv}
  \item \textbf{WHO GHO} --- Antenatal care coverage (4+ visits):\\
  \url{https://apps.who.int/gho/athena/api/GHO/WHS4_154?format=csv}
  \item \textbf{World Bank} --- Female literacy rate, GDP per capita, health expenditure:\\
  \url{https://data.worldbank.org/}
  \item \textbf{DHS Program} (Demographic and Health Surveys) --- survey-level indicators:\\
  \url{https://www.statcompiler.com/}\\
  Provides subnational data on maternal health, contraceptive use, and antenatal care for African and Asian countries.
\end{itemize}

\textbf{Required analyses.}
\begin{enumerate}
  \item Load MMR data and filter to the most recent year. Identify the 20 countries with the highest MMR.
  \item Merge with skilled birth attendance, antenatal care coverage, and socioeconomic indicators.
  \item Exploratory analysis: scatter plots of MMR vs.\ each predictor. Compute Pearson and Spearman correlations.
  \item Multiple linear regression: which factors significantly predict MMR? Check model assumptions (residual plots, normality).
  \item Create a choropleth map of MMR across Africa.
  \item (Bonus) Compare MMR trends over time (2000--2020) for 5 countries that implemented specific maternal health programs.
\end{enumerate}

\textbf{Deliverables.}
\begin{itemize}
  \item Scatter plots with regression lines (MMR vs.\ skilled birth attendance, MMR vs.\ health expenditure).
  \item Regression table with coefficients, standard errors, p-values, and $R^2$.
  \item Choropleth map of maternal mortality in Africa.
  \item Discussion of policy implications: which interventions could reduce MMR?
\end{itemize}

\begin{clinicalnote}
Maternal mortality is one of the starkest health inequities in the world. Sub-Saharan Africa accounts for approximately 70\% of global maternal deaths. Most of these deaths are preventable with skilled birth attendants, emergency obstetric care, and adequate antenatal visits.
\end{clinicalnote}

% ------------------------------------------------------------
\subsection{Project 4: Heart disease prediction model comparison}

\textbf{Objective.} Compare multiple machine learning models for predicting heart disease and identify the most important clinical predictors.

\textbf{Dataset.}
\begin{itemize}
  \item \textbf{UCI Heart Disease Dataset} (Cleveland):\\
  \url{https://archive.ics.uci.edu/ml/machine-learning-databases/heart-disease/processed.cleveland.data}\\
  303 patients, 13 clinical features (age, sex, chest pain type, resting BP, cholesterol, fasting blood sugar, resting ECG, max heart rate, exercise-induced angina, ST depression, slope, number of major vessels, thalassemia), binary target (presence of heart disease).
  \item Column names: \texttt{age, sex, cp, trestbps, chol, fbs, restecg, thalach, exang, oldpeak, slope, ca, thal, target}.
\end{itemize}

\textbf{Required analyses.}
\begin{enumerate}
  \item Load and clean the data. Handle missing values (coded as \texttt{?} in the original file).
  \item Exploratory analysis: distributions of each feature by heart disease status.
  \item Train-test split (80/20). Use stratified splitting to maintain class balance.
  \item Train and evaluate 4 models:
  \begin{itemize}
    \item Logistic regression
    \item Decision tree
    \item Random forest
    \item K-nearest neighbors (KNN)
  \end{itemize}
  \item For each model: report accuracy, precision, recall, F1-score, and AUC-ROC.
  \item Plot ROC curves for all 4 models on the same figure.
  \item Feature importance: which clinical variables are the strongest predictors?
\end{enumerate}

\textbf{Deliverables.}
\begin{itemize}
  \item Comparison table of model performance metrics.
  \item ROC curves for all models on one plot.
  \item Feature importance bar chart (from random forest).
  \item Confusion matrix for the best-performing model.
  \item Clinical interpretation: do the top predictors align with cardiology guidelines?
\end{itemize}

\begin{pythontip}
Use \texttt{sklearn.model\_selection.cross\_val\_score} with 5-fold cross-validation instead of a single train-test split. This gives a more reliable estimate of model performance, especially with a small dataset like this one (303 patients).
\end{pythontip}

% ------------------------------------------------------------
\subsection{Project 5: Malaria burden and intervention effectiveness}

\textbf{Objective.} Analyze trends in malaria burden across African countries and evaluate whether key interventions (bed nets, indoor spraying, treatment) are associated with declining incidence.

\textbf{Datasets.}
\begin{itemize}
  \item \textbf{WHO GHO} --- Malaria incidence (per 1\,000 population at risk):\\
  \url{https://apps.who.int/gho/athena/api/GHO/MALARIA_INCIDENCE?format=csv}
  \item \textbf{WHO GHO} --- Malaria mortality rate:\\
  \url{https://apps.who.int/gho/athena/api/GHO/MALARIA_DEATHS_PER100000?format=csv}
  \item \textbf{WHO World Malaria Report data} --- intervention coverage (ITN use, IRS coverage, ACT treatment):\\
  \url{https://www.who.int/teams/global-malaria-programme/reports/world-malaria-report}\\
  Annex data tables are available as Excel downloads.
  \item \textbf{The Malaria Atlas Project}:\\
  \url{https://malariaatlas.org/}\\
  Subnational prevalence estimates and raster maps.
  \item \textbf{DHS Program}:\\
  \url{https://www.statcompiler.com/}\\
  Survey-based data on ITN ownership and use, IRS coverage, and treatment-seeking behavior.
\end{itemize}

\textbf{Required analyses.}
\begin{enumerate}
  \item Load malaria incidence and mortality data. Filter to sub-Saharan African countries.
  \item Plot incidence trends (2000--2022) for 10 countries with the highest burden.
  \item Merge with intervention coverage data (ITN use, IRS, ACT access).
  \item Compute correlations between changes in intervention coverage and changes in incidence over time (i.e., does increasing ITN use predict decreasing incidence?).
  \item Create a choropleth map of current malaria incidence across Africa.
  \item Linear regression: predict change in malaria incidence from change in intervention coverage, controlling for GDP per capita and urbanization.
  \item (Bonus) Create a before-and-after comparison for countries that scaled up ITN distribution after the Global Fund grants.
\end{enumerate}

\textbf{Deliverables.}
\begin{itemize}
  \item Time series plot of malaria incidence for 10 countries (2000--2022).
  \item Scatter plot of ITN use change vs.\ incidence change with regression line.
  \item Choropleth map of malaria incidence across Africa.
  \item Regression table with interpretation.
  \item Policy discussion: which interventions show the strongest association with declining malaria?
\end{itemize}

\begin{clinicalnote}
Between 2000 and 2015, malaria mortality in Africa declined by approximately 44\%, largely driven by the scale-up of ITNs, IRS, and artemisinin-based combination therapies (ACTs). Since 2015, progress has stalled, making this analysis timely and policy-relevant.
\end{clinicalnote}

% ============================================================
\section{Report template}

Your written report should follow this structure (5--8 pages, excluding code appendix):

\subsection{1. Introduction (0.5--1 page)}

\begin{itemize}
  \item Health context: why does this question matter?
  \item Specific research question or objective.
  \item Brief description of the dataset(s) used.
\end{itemize}

\subsection{2. Data description (1 page)}

\begin{itemize}
  \item Source(s) and citation(s).
  \item Number of observations and variables.
  \item Key variables: name, type, unit, and clinical meaning.
  \item Summary statistics table (\texttt{.describe()}).
  \item Missing data: how much, which variables, how handled.
\end{itemize}

\subsection{3. Methods (1 page)}

\begin{itemize}
  \item Data cleaning steps (handling missing values, recoding variables, outlier treatment).
  \item Statistical tests used and why (e.g., ``we used the chi-squared test because both variables are categorical'').
  \item Models used (logistic regression, random forest, etc.) with hyperparameters.
  \item Validation strategy (train-test split, cross-validation).
\end{itemize}

\subsection{4. Results (1.5--2 pages)}

\begin{itemize}
  \item Figures and tables with captions.
  \item Report test statistics, p-values, confidence intervals.
  \item Model performance metrics.
  \item Do not interpret here --- just present the numbers.
\end{itemize}

\subsection{5. Discussion (1--1.5 pages)}

\begin{itemize}
  \item What do the results mean clinically?
  \item How do your findings compare with published literature?
  \item What surprised you?
  \item What are the practical implications?
\end{itemize}

\subsection{6. Limitations (0.5 page)}

\begin{itemize}
  \item Data quality issues (missing data, measurement error, selection bias).
  \item Generalizability: does the dataset represent the population you care about?
  \item What analyses would you do with more time or data?
\end{itemize}

\begin{warningbox}
Every figure must have a title, labeled axes, and a caption explaining what the reader should see. Every table must have a header row and a caption. Figures without labels are the most common reason for lost marks.
\end{warningbox}

% ============================================================
\section{Presentation guidelines}

You will give a 5-minute oral presentation followed by 2--3 minutes of questions.

\subsection{Structure}

\begin{enumerate}
  \item \textbf{Slide 1: Title.} Project title, your name, date.
  \item \textbf{Slide 2: Motivation.} Why does this health question matter? One or two key facts.
  \item \textbf{Slide 3: Data.} What data did you use? How many observations? What are the key variables?
  \item \textbf{Slide 4--5: Key results.} Your 2--3 most important figures or tables. Walk the audience through each one.
  \item \textbf{Slide 6: Conclusions.} One or two take-home messages. What should the audience remember?
\end{enumerate}

\subsection{Tips}

\begin{itemize}
  \item \textbf{5 minutes is short.} Practice with a timer. Cut anything that is not essential.
  \item \textbf{One message per slide.} Do not crowd slides with text.
  \item \textbf{Speak to the audience, not the screen.}
  \item \textbf{Anticipate questions.} Know your limitations and be ready to discuss them.
  \item \textbf{No code on slides.} Show results, not the code that produced them.
\end{itemize}

\begin{pythontip}
Save your figures as high-resolution PNGs for slides:
\begin{minted}{python}
fig.savefig("figure1.png", dpi=300, bbox_inches="tight")
\end{minted}
This avoids blurry screenshots and ensures your plots look professional in a presentation.
\end{pythontip}

% ============================================================
\section{Grading rubric}

The capstone project is graded out of 100 points:

\begin{center}
\begin{tabular}{lcp{8cm}}
\toprule
\textbf{Component} & \textbf{Points} & \textbf{Criteria} \\
\midrule
Research question & 10 & Clear, specific, and health-relevant \\
Data handling & 15 & Proper loading, cleaning, missing value treatment, and documentation \\
Exploratory analysis & 15 & Appropriate summary statistics, distributions, and at least 3 well-labeled visualizations \\
Statistical analysis & 20 & Correct choice and application of tests or models; proper interpretation of p-values, confidence intervals, and effect sizes \\
Report quality & 20 & Follows the template; clear writing; figures have titles, labels, and captions; limitations honestly discussed \\
Code quality & 10 & Reproducible notebook that runs top-to-bottom; comments explain key steps; no unnecessary code \\
Oral presentation & 10 & Clear, within time limit, answers questions confidently \\
\midrule
\textbf{Total} & \textbf{100} & \\
\bottomrule
\end{tabular}
\end{center}

\subsection{Grade boundaries}

\begin{itemize}
  \item \textbf{A (90--100):} Excellent analysis with insightful interpretation. Code is clean and well-documented. Goes beyond the minimum requirements.
  \item \textbf{B (75--89):} Solid analysis with correct methodology. Minor issues in presentation or interpretation.
  \item \textbf{C (60--74):} Meets basic requirements but has notable gaps (e.g., missing visualizations, superficial interpretation, code errors).
  \item \textbf{D/F ($<60$):} Incomplete analysis, major methodological errors, or non-reproducible code.
\end{itemize}

% ============================================================
\section{Getting started: a checklist}

\begin{exercise}
Before you leave today, complete this checklist:
\begin{enumerate}
  \item Choose your project (1--5 or a custom proposal).
  \item Download your primary dataset. Load it into a Jupyter notebook and run \texttt{.head()}, \texttt{.shape}, \texttt{.info()}, and \texttt{.describe()}.
  \item Write your research question in one sentence.
  \item List 3 visualizations you plan to create.
  \item List 1 statistical test or model you plan to use.
  \item Identify potential problems: missing data, messy column names, multiple files to merge.
  \item Create a project folder with the structure:
\begin{minted}{text}
capstone_project/
    data/
        raw/          # original downloaded files
        cleaned/      # processed data
    notebooks/
        01_exploration.ipynb
        02_analysis.ipynb
    figures/
    report.pdf
\end{minted}
  \item Submit your project choice and research question to the instructor by the end of the week.
\end{enumerate}
\end{exercise}

% ============================================================
\section{Chapter summary}

\begin{itemize}
  \item The capstone project integrates all course skills: data loading, cleaning, exploration, visualization, statistical analysis, and interpretation.
  \item Five suggested projects cover diabetes, COVID-19, maternal mortality, heart disease, and malaria --- all using real, publicly available datasets.
  \item Each project requires a clear health question, reproducible code, at least 3 visualizations, and at least 1 statistical test or model.
  \item The report follows a standard structure: Introduction, Data, Methods, Results, Discussion, Limitations.
  \item Presentations are 5 minutes --- focus on motivation, key results, and conclusions.
  \item The grading rubric rewards clear thinking and honest interpretation over complex methodology.
\end{itemize}
